\chapter{Critical-line symmetry and the arithmetic target}

An involution can preserve residue data without making the dilation
operator symmetric. This distinction is visible in the signs of
full principal parts.

\section{Inversion and conjugation}

For real $c$, define on $\mathscr H_c$
\[
 (J_cf)(x)=x^{-2c}\overline{f(x^{-1})}.
\]
In logarithmic coordinates,
$U_cJ_cf(t)=\overline{U_cf(-t)}$. Consequently $J_c$ is an
antiunitary involution. Direct substitution in the Mellin integral
gives
\begin{equation}
 \mathcal M(J_cf)(s)=\overline{F(2c-\overline s)},
 \qquad J_cP_c=P_cJ_c
 \label{mo:inversion}
\end{equation}
on compact smooth tests, and then on the closed operator domain.
For the domain assertion, logarithmic reflection and conjugation
preserve the weak-derivative norm from
Lemma~\ref{mo:mellin-unitary}.

Suppose now that the entire function satisfies
\begin{equation}
 \Xi(2c-\overline s)=\overline{\Xi(s)}.
 \label{mo:entire-symmetry}
\end{equation}
Write $\rho^\#=2c-\overline\rho$. Its zeros and multiplicities
are preserved by $\rho\mapsto\rho^\#$. Choose the finite zero set
$Z$ invariant under this map.

\begin{proposition}[the residue involution]
\label{mo:residue-involution}
On full jet coordinates, inversion acts by
\[
 (J_Za)_{\rho,k}=(-1)^{k+1}\overline{a_{\rho^\#,k}}.
\]
Thus $\mathcal R_ZJ_c=J_Z\mathcal R_Z$, and $J_Z$ is an
antiunitary involution for the Euclidean jet norm. It commutes
with $K_c$. At a simple critical-line zero, the residue action
is minus complex conjugation.
\end{proposition}
\begin{proof}
Put $s=\rho+z$. Equations \eqref{mo:inversion} and
\eqref{mo:entire-symmetry} identify the quotient for $J_cf$ with
the conjugate of the quotient for $f$ at
$\rho^\#-\overline z$. A coefficient of order $-(k+1)$ therefore
acquires the factor $(-1)^{k+1}$. This proves the formula.
The norm and involution assertions follow by permuting coordinates
and conjugating them. For commutation, note that
$z_{\rho^\#,c}=\overline{z_{\rho,c}}$ and that consecutive signs
in the formula are opposite. Substitution into
\eqref{mo:jet-generator} gives $J_ZK_c=K_cJ_Z$.
\end{proof}

The same assertion holds for positive diagonal jet weights invariant
under $\rho\mapsto\rho^\#$. Commutation with this antiunitary does
not remove a Jordan block. The symmetry criterion of
Theorem~\ref{mo:finite-symmetry} still requires simple zeros on the line.

The involution extends to compactly supported distributions by
\[
 \langle J_cu,\phi\rangle
 =\overline{\left\langle u,
         x^{2c-2}\overline{\phi(x^{-1})}\right\rangle}.
\]
Changing variables proves agreement with its action on functions.
For occupation distributions,
\[
 J_cTv_m=m^{1-2c}\delta_m.
\]
At $c=1/2$ this is $\delta_m$, with Mellin transform $m^{s-1}$.
Its $Q$ eigenvalue is $-\log m$. Except for $m=1$, it lies outside
the distributional image supported at the reciprocal integers.
Thus an inversion-invariant space must contain both signs of this
logarithmic position observable. Identifying this involution with
an arithmetic modular conjugation requires a separate map of the
actual spaces and pairings.

\section{The completed Riemann function}

The preceding results apply to any nonzero entire $\Xi$.
For the arithmetic specialization, an integral formula fixes the
normalization and both symmetries without a choice of zero locations.
For $t>0$, put
\[
 \vartheta(t)=\sum_{n\in\mathbb Z}e^{-\pi n^2t},\qquad
 \psi(t)=\sum_{n\geq1}e^{-\pi n^2t}.
\]
These series and their derivatives converge uniformly on compact
subsets of $t>0$.

The Gaussian identity
$\int_{\mathbb R}e^{-\pi t x^2}e^{-2\pi i yx}\,dx
=t^{-1/2}e^{-\pi y^2/t}$ follows by differentiation in $y$ and
integration by parts; its value at zero is the Gaussian integral.
Periodize $e^{-\pi t x^2}$ over the integers. Its Fourier coefficients
are the displayed Gaussian values at integer $y$. Both the
periodized series and the resulting Fourier series converge
absolutely and uniformly. Equality follows from uniqueness of
Fourier coefficients for continuous periodic functions. For
completeness, uniqueness follows by convolution with
\[
 F_n(x)=\frac1n\left|\sum_{k=0}^{n-1}e^{2\pi i kx}\right|^2:
\]
these nonnegative kernels have integral one and their mass outside
any neighborhood of the integers tends to zero. Their convolutions
converge uniformly to a continuous periodic function. Evaluating
at zero gives
\[
 \vartheta(t)=t^{-1/2}\vartheta(t^{-1}),\qquad
 \psi(t)=t^{-1/2}\psi(t^{-1})+\tfrac12(t^{-1/2}-1).
\]

Define
\begin{equation}
 \xi_{\mathbb R}(s)=\frac12+
 \frac{s(s-1)}2\int_1^\infty
 \psi(t)\bigl(t^{s/2}+t^{(1-s)/2}\bigr)\,\frac{dt}{t}.
 \label{mo:xi-integral}
\end{equation}
The exponential decrease of $\psi$ makes this integral entire,
with every derivative locally uniformly convergent. It gives
$\xi_{\mathbb R}(0)=\xi_{\mathbb R}(1)=1/2$ and
\[
 \xi_{\mathbb R}(1-s)=\xi_{\mathbb R}(s),\qquad
 \xi_{\mathbb R}(\overline s)=\overline{\xi_{\mathbb R}(s)}.
\]
In particular it satisfies \eqref{mo:entire-symmetry} at $c=1/2$.

On $\operatorname{Re}s>1$, define
$\zeta(s)=\sum_{n\geq1}n^{-s}$ and
$\Gamma(s/2)=\int_0^\infty e^{-u}u^{s/2-1}\,du$.
Absolute convergence permits summation under the integral and gives
\[
 \int_0^\infty\psi(t)t^{s/2}\,\frac{dt}{t}
       =\pi^{-s/2}\Gamma(s/2)\zeta(s).
\]
Split the integral at one and use the theta identity on $(0,1)$.
The elementary contribution is
$1/(s-1)-1/s=1/(s(s-1))$; the remaining integral is that in
\eqref{mo:xi-integral}. Therefore
\[
 \xi_{\mathbb R}(s)=
 \tfrac12s(s-1)\pi^{-s/2}\Gamma(s/2)\zeta(s)
 \quad(\operatorname{Re}s>1).
\]
Equation \eqref{mo:xi-integral} supplies the entire continuation.
No assertion about simplicity or critical-line placement of its
zeros was used.

\begin{proposition}[the occupation vacuum and the zero divisor]
\label{mo:vacuum}
One has $1/4<\xi_{\mathbb R}(1/2)<1/2$.
Suppose a self-adjoint operator $H$ has the property that every
eigenvalue of $1/2+iH$ is a zero of $\xi_{\mathbb R}$.
There is no injective linear map $A$ whose domain contains $v_1$,
whose image of $v_1$ belongs to $\operatorname{Dom}H$, and which
satisfies $ANv_1=HAv_1$.
\end{proposition}
\begin{proof}
At $s=1/2$, \eqref{mo:xi-integral} reads
\[
 \xi_{\mathbb R}(1/2)=\frac12-
          \frac14\int_1^\infty\psi(t)t^{-3/4}\,dt.
\]
For $t\geq1$, the inequality $n^2\geq n$ gives
$0<\psi(t)\leq e^{-\pi t}/(1-e^{-\pi})$.
Consequently the positive integral is bounded above by
$e^{-\pi}/[\pi(1-e^{-\pi})]<1$. This proves the bounds.
Since $Nv_1=0$, the proposed identity gives $HAv_1=0$.
Injectivity makes $Av_1\ne0$, so $1/2$ would be an eigenvalue
of $1/2+iH$ and a zero of $\xi_{\mathbb R}$. The bounds exclude it.
\end{proof}

For $\Xi=\xi_{\mathbb R}$ and $c=1/2$, the full jet eigenvalue is
\[
 z_{\rho,1/2}=-i(\rho-1/2)
       =\operatorname{Im}\rho+i(1/2-\operatorname{Re}\rho).
\]
The affine shift $1/2+iK_{1/2}$ has eigenvalues $\rho$ and, on
full jets, superdiagonal $1$. Replacing $z_{\rho,1/2}$ by
$\operatorname{Im}\rho$ projects the eigenvalue to the critical
line. Removing the superdiagonal removes the generalized-eigenvector
data at multiple zeros. Both changes require an explicit comparison.

If the selected zeros are simple and on the critical line, their
weighted evaluation completion from
Theorem~\ref{mo:countable-completion} has a self-adjoint diagonal dilation
operator. Its spectrum and norm were specified through the selected
zero data and weights. This conditional construction does not
produce an independently defined arithmetic positive pairing or
establish critical-line placement. With full multiple-zero jets,
Theorem~\ref{mo:finite-symmetry} gives an additional obstruction to a
positive symmetric realization.

\section{Local quotient representations and traces}

Let $\mathcal O_\rho$ be the algebra of holomorphic germs at
$\rho$. If $\Xi$ has a zero of order $r$ there, the linear map
\[
 \mathcal O_\rho/(\Xi)\longrightarrow\mathbb C^r,
 \qquad [F]\longmapsto(R_{\rho,k}F)_{0\leq k<r}
\]
is an isomorphism. Indeed $F/\Xi$ has no principal part exactly
when $F$ is divisible by $\Xi$ as a germ. Taylor reduction modulo
$(s-\rho)^r$ and \eqref{mo:jet-formula} prove surjectivity.
Multiplication by $s$ becomes $\rho I+S$, where
$(Sa)_k=a_{k+1}$. Thus the affine operator $c+iK_c$ is exactly
the local multiplication operator. This comparison preserves
multiplicity and its nilpotent action.

This local quotient is the algebraic form of zero multiplicity.
Meyer constructs a global virtual representation of the idele
class group on nuclear bornological spaces. Its character is
the Weil distribution, and its spectral multiplicities are the
pole orders of the completed $L$-function
\cite[Theorems~5.8 and~5.11, pp.~43 and~47]{Meyer}.
Zeros contribute through the negative representation. That result
does not specify the positive evaluation weights constructed here.
An identification with these Hilbert spaces requires maps from
the global representation, continuity in both norms, and compatibility
with the scaling action and the chosen pairing.

For finite blocks the distinction between trace and representation
is explicit. The trace $\operatorname{tr}p(\rho I+S)=rp(\rho)$
equals the trace on $r$ scalar copies of $\rho$.
When $r>1$ these representations are not isomorphic: the first has
a nonzero nilpotent part, while the second has none.
For infinite blocks an ordinary operator trace additionally requires
trace-class convergence. In the simple critical-line case,
\[
 \operatorname{Tr}\varphi(K_w)=\sum_j\varphi(\gamma_j)
 \quad\text{if}\quad \sum_j|\varphi(\gamma_j)|<\infty.
\]
To prove this formula, use the orthonormal basis
$w_j^{-1/2}e_j$ and the diagonal trace-class criterion.
In particular the weights disappear from the trace.
Hence this trace cannot recover the positive form chosen on
the test functions.

The occupation operator has a different trace. For real $\beta>1$,
\[
 \operatorname{Tr}e^{-\beta N}=\sum_{m\geq1}m^{-\beta}
 =\zeta(\beta).
\]
The positive diagonal sum converges by integral comparison and
diverges for $\beta\leq1$. A trace over zero ordinates involves
a different index set and requires a proved comparison of the
operators and trace functionals. Mellin evaluation and equality
of local multiplicities do not supply that comparison.

\section{The occupation realization that remains}

An occupation-to-zero realization would require a Hilbert space
$\mathscr H_\xi$, an independently specified positive inner product,
and a self-adjoint operator $H$ there. It would also require an
injective linear map $A$ on a dense $N$-invariant domain
$\mathscr D_N\subset\operatorname{Dom}N$, with dense range in
$\mathscr H_\xi$, such that
\[
 A\mathscr D_N\subset\operatorname{Dom}H,\qquad
 AN=HA\quad\text{on }\mathscr D_N.
\]
Closedness and domain compatibility are part of this operator
problem. Matching the arithmetic zero divisor further requires
the eigenvalues of $1/2+iH$, with their multiplicities, to give
the actual zeros of $\xi_{\mathbb R}$. Any determinant or trace
comparison requires its own definition, convergence domain, and
normalization.

The explicit map $T$ constructs the occupation-to-distribution
edge on every finite sequence. Its residue composition fails the
required energy identity by Proposition~\ref{mo:occupation-defect}, including
multiple zeros. A domain containing the finite occupation core
cannot fix that algebraic failure by completion. Domains excluding
that core are different operator problems and require new maps.
Proposition~\ref{mo:vacuum} also excludes every injective map
with the actual zero-divisor target when its domain contains $v_1$.
The normalized dilation operator instead has the exact jet
comparison Proposition~\ref{mo:jet-comparison}; it acts on a continuous
spectral carrier. Its evaluation map is not a closable transfer
from the original Hilbert space in the settings proved above.

The physical distinction is between logarithmic position and its
conjugate momentum. Integer occupation selects discrete position
values. Mellin dilation measures continuous momentum. Residue
evaluation samples a function of that momentum at a chosen divisor;
a positive sampling norm changes the state space. The displayed
maps and obstructions specify these operations. They supply no
identification of their thermal traces or of an arithmetic trace
with the sampled norm.

\begin{thebibliography}{9}
\bibitem{Meyer}
R.~Meyer, \emph{On a representation of the idele class group related
to primes and zeros of $L$-functions}, Duke Math. J. \textbf{127}
(2005), 519--595. \href{https://arxiv.org/abs/math/0311468v3}
{arXiv:math/0311468v3}. Theorem~5.8 gives the character formula;
Theorem~5.11 gives the spectral multiplicities.
\end{thebibliography}
